\section{Tessuto nervoso}

\textbf{Tessuto nervoso}: mette in comunicazione le varie parti dell'organismo, riceve informazioni dall'ambiente esterno o interno e risponde alle informazioni ricevute.

\rubrica{Suddivisione anatomica}
\begin{itemize}
  \item \textbf{sistema nervoso centrale (SNC)}: encefalo (nella scatola cranica) e midollo spinale (nel rachide);
  \item \textbf{sistema nervoso periferico (SNP)}.
\end{itemize}

\rubrica{Tipi di cellule}
\begin{itemize}
  \item \textbf{cellule nervose} o \textbf{neuroni};
  \item \textbf{cellule di sostegno} o \textbf{neuroglia} (glia).
\end{itemize}

% ---------------------------------------------------------------------------
\subsection{Il neurone}

\begin{itemize}
  \item \textbf{corpo cellulare}: sito del normale metabolismo cellulare;
  \item \textbf{dendriti}: prolungamenti cellulari che ricevono il potenziale d'azione e lo conducono verso il corpo cellulare (\textit{polo ricevente});
  \item \textbf{assone} (o neurite): prolungamento cellulare che conduce il potenziale d'azione lontano dal corpo cellulare (\textit{polo trasmittente}).
\end{itemize}

\rubrica{Corpo cellulare (soma, pirenoforo)}
\begin{itemize}
  \item un solo nucleo, con nucleolo evidente;
  \item RER molto esteso;
  \item numerosi apparati del Golgi;
  \item mitocondri e altri organuli;
  \item numerosi filamenti intermedi (\textbf{neurofilamenti});
  \item i microtubuli separano il RER in ampie aree (\textbf{sostanza tigroide} o \textbf{corpi di Nissl}).
\end{itemize}

\rubrica{Organuli del corpo cellulare}
\begin{itemize}
  \item \textbf{membrana plasmatica}: sede dei fenomeni bioelettrici;
  \item \textbf{nucleo}: unico, voluminoso, tondeggiante, al centro del pirenoforo, con uno o più nucleoli;
  \item \textbf{reticolo endoplasmatico liscio}: presente nei dendriti e nell'assone; regolazione del calcio (alti livelli nei dendriti);
  \item \textbf{apparato di Golgi}: molto sviluppato, con molte vescicole associate, in continuità con il REL;
  \item \textbf{mitocondri}: abbondanti nel corpo cellulare e nei bottoni presinaptici;
  \item \textbf{lisosomi}: nel soma, in regione perinucleare; presenti inclusi citoplasmatici e granuli di lipofuscina;
  \item \textbf{centrioli}: una sola coppia, sede dei centri di organizzazione microtubulare.
\end{itemize}

\rubrica{Citoscheletro}
Particolarmente sviluppato.
\begin{itemize}
  \item \textbf{neurofilamenti}: filamenti intermedi ($\phi$ 8--10\,nm), presenti nel pirenoforo, nei dendriti e nell'assone; più abbondanti nell'assone;
  \item \textbf{microtubuli}: $\phi$ 20--30\,nm ($\phi$ del lume 15\,nm), presenti uniformemente in tutta la cellula; più abbondanti nei dendriti;
  \item \textbf{microfilamenti di actina}: $\phi$ 5--6\,nm, associati al plasmalemma.
\end{itemize}

\rubrica{Dendriti}
\begin{itemize}
  \item prolungamenti citoplasmatici rivestiti da plasmalemma; più grossi vicino al pirenoforo, più sottili nelle suddivisioni;
  \item relativamente corti (max 700\,$\mu$m);
  \item mitocondri orientati longitudinalmente;
  \item \textbf{spine dendritiche}: estroflessioni di forma allungata con rigonfiamento apicale, sede di sinapsi assospinodendritiche;
  \item sinapsi assodendritiche e numerose sinapsi;
  \item conferiscono al neurone la capacità di integrare informazioni provenienti da più fonti $\rightarrow$ rappresentano l'apparato ricevente del neurone.
\end{itemize}

\rubrica{Assone}
\begin{itemize}
  \item un solo assone origina dal corpo cellulare $\rightarrow$ \textbf{cono di emergenza};
  \item al cono di emergenza segue il \textbf{segmento iniziale}; da queste due regioni originano i potenziali d'azione trasportati lungo la fibra nervosa;
  \item ramificazioni secondarie, se presenti, ad angolo retto;
  \item diametro costante compreso tra 0,1--20\,$\mu$m;
  \item lunghezza da pochi micron fino a più di 1 metro;
  \item contiene mitocondri, microtubuli, neurofilamenti, microfilamenti e vescicole disposti parallelamente all'asse maggiore.
\end{itemize}

\rubrica{Flusso assonico}
Trasporto garantito nei due sensi da motori molecolari lungo i microtubuli:
\begin{itemize}
  \item \textbf{trasporto anterogrado} (rapido): mediato dalla \textbf{chinesina}, verso l'estremità positiva del microtubulo (terminazione sinaptica);
  \item \textbf{trasporto retrogrado}: mediato dalla \textbf{dineina}, verso l'estremità negativa (corpo cellulare).
\end{itemize}

% ---------------------------------------------------------------------------
\subsection{Classificazione strutturale dei neuroni}

\begin{itemize}
  \item \textbf{unipolari}: un solo prolungamento;
  \item \textbf{bipolari}: un dendrite e un assone ben distinti, separati dal corpo cellulare da cui fuoriescono da due posizioni opposte;
  \item \textbf{pseudounipolari}: i due tipi di processi originano da un unico polo e si separano a distanza dal corpo cellulare;
  \item \textbf{multipolari}: molti dendriti e un solo assone (a questa categoria appartengono i \textbf{neuroni di Purkinje}).
\end{itemize}

% ---------------------------------------------------------------------------
\subsection{Neuroglia del SNC}

\rubrica{Astrociti}
\begin{itemize}
  \item grandi dimensioni, forma stellata;
  \item costituiscono manicotti attorno ai vasi sanguigni dell'encefalo (\textbf{pedicelli terminali});
  \item mantengono la \textbf{barriera ematoencefalica}, promuovendo la formazione di giunzioni comunicanti tra le cellule endoteliali dei capillari;
  \item impalcatura tridimensionale di sostegno per i neuroni;
  \item provvedono alla riparazione del tessuto nervoso danneggiato.
\end{itemize}

\rubrica{Oligodendrociti}
\begin{itemize}
  \item rari processi dendritici, cellule stellate;
  \item formano la \textbf{guaina mielinica}, aumentando la velocità dello stimolo nervoso;
  \item ogni oligodendrocita è capace di avvolgere più assoni.
\end{itemize}

\rubrica{Cellule ependimali}
\begin{itemize}
  \item cellule cubiche o cilindriche che rivestono le cavità del SNC (ventricoli e canale centrale), con produzione di \textbf{liquor cefalorachidiano};
  \item \textbf{taniciti}: cellule ependimali atipiche in cui la parte basale si espande perpendicolarmente alla parete ventricolare, approfondendosi nel tessuto nervoso.
\end{itemize}

\rubrica{Microglia}
\begin{itemize}
  \item piccole dimensioni;
  \item \textbf{macrofagi specializzati del SNC};
  \item cellule di difesa mobili di tipo fagocitico, per rimuovere detriti cellulari, rifiuti ed agenti patogeni;
  \item in seguito a infiammazioni, traumi o ischemie, vengono attivate (da \textit{microglia ramificata} a \textit{microglia attivata}) e migrano verso le aree interessate dal danno.
\end{itemize}

% ---------------------------------------------------------------------------
\subsection{Neuroglia del SNP}

\rubrica{Cellule di Schwann (neurilemmociti)}
\begin{itemize}
  \item presenti nel SNP;
  \item cellule appiattite, con ampio citoplasma ricco di mielina;
  \item responsabili della mielinizzazione degli assoni periferici;
  \item formano la guaina mielinica attorno ad un solo assone (arrotolamento a spirale a partire dal \textbf{mesassone}).
\end{itemize}

\rubrica{Cellule satelliti (amficiti)}
\begin{itemize}
  \item presenti nel SNP;
  \item circondano i corpi cellulari dei neuroni dei gangli;
  \item regolano i livelli di $O_2$, $CO_2$, nutrienti e neurotrasmettitori.
\end{itemize}

% ---------------------------------------------------------------------------
\subsection{Sostanza grigia e sostanza bianca}

\begin{itemize}
  \item \textbf{sostanza grigia} (scarsamente mielinizzata): vi sono concentrati i corpi cellulari, i dendriti, i terminali degli assoni e vari tipi di cellule gliali;
  \item \textbf{sostanza bianca} (mielinizzata): vi sono concentrati gli assoni coperti dalla guaina mielinica -- formata nel SNC dagli oligodendrociti e nel SNP dalle cellule di Schwann -- e vari tipi di cellule gliali.
\end{itemize}

% ---------------------------------------------------------------------------
\subsection{Fibra nervosa}

Costituita dall'assone ricoperto dal proprio involucro formato da cellule gliali.
\begin{itemize}
  \item \textbf{fibra mielinica}: la cellula di Schwann si arrotola a spirale attorno all'assone, circondandolo con più giri di membrana (\textbf{guaina mielinica}) e, in superficie, con citoplasma e membrana plasmatica (\textbf{neurilemma}); la guaina è interrotta a intervalli dai \textbf{nodi di Ranvier};
  \item \textbf{fibra amielinica}: una sola cellula di Schwann ricopre più assoni con un solo giro di membrana e citoplasma.
\end{itemize}

% ---------------------------------------------------------------------------
\subsection{Nervo periferico}

Involucri connettivali, dall'esterno verso l'interno:
\begin{itemize}
  \item \textbf{epinevrio}: copre il nervo periferico;
  \item \textbf{perinevrio}: avvolge ogni fascicolo;
  \item \textbf{endonevrio}: attorno alla singola fibra.
\end{itemize}
Nel nervo decorrono vasi sanguigni; ogni fascicolo raccoglie fibre costituite dall'assone mielinico e dalla cellula di Schwann.
